\documentclass[10pt,a4paper]{amsart}
\usepackage[margin=19mm]{geometry}
\usepackage{amsmath,amssymb,mathtools}
\usepackage[scr=rsfs,cal=euler]{mathalfa}
\usepackage{xfrac}
\usepackage[unicode,hidelinks,bookmarks=false]{hyperref}
\newcommand{\ie}{i.e.\ }
\newcommand{\cf}{cf.\ }
\newcommand{\resp}{respectively\ }
\newcommand{\Subsection}{\subsection}
\providecommand{\nobreakdash}{\nobreak\hbox{-}}
\setlength{\emergencystretch}{2em}
\multlinegap0pt
\usepackage{float}
\usepackage{mathtools}
\usepackage{amssymb}
\usepackage{booktabs}
\usepackage{xcolor}
\newtheorem{theorem}{Theorem}
\title{Chaos Analysis in the Hybrid Quintic Duffing-Riemann Zeta System via Decomposition\\}
\author{Zeraoulia Rafik, Pedro Caceres}
\date{Source: arXiv:2212.12438v2 (CC BY 4.0)}
\begin{document}
\maketitle

\noindent\emph{Source and licence.} Chaos Analysis in the Hybrid Quintic Duffing-Riemann Zeta System via Decomposition\\. Version arXiv:2212.12438v2 (CC BY 4.0), distributed by arXiv under the Creative Commons Attribution 4.0 International licence, http://creativecommons.org/licenses/by/4.0/, as recorded in the arXiv licence metadata for this exact version and shown on the abstract page https://arxiv.org/abs/2212.12438v2. Full licence notice: \texttt{licenses/arxiv-2212.12438-CC-BY-4.0.txt}.\par\smallskip

\noindent\textit{Editorial scope.} Counted unit: the article's theorem thm:chaos\_control (source0.tex lines 1469--1476). The article's complete original proof of it is delivered with it (source0.tex lines 1478--1484, one \texttt{proof} environment of 7 lines). No mathematics is written, repaired or re-derived: the statement and the proof are the article's own lines, in the article's own notation, and no retained line is substituted or re-flowed.  No further source-local context is needed: every cross-reference of every retained line resolves inside the two delivered spans.  Nothing was omitted from any retained span, so the package carries no omission record.  No retained line contains a citation, so the wrapper carries no bibliography and no bibliography entry.  No retained line contains a figure environment, an image-inclusion command or drawing code, so no image is delivered and the package declares no asset dependency.  Editorial formatting only, in addition: an AMS \texttt{amsart} preamble that restates the article's own declarations and macros used by the retained lines, namely the environments \texttt{theorem} on separate counters, together with the standard packages those lines need.  The retained environments print in this document as Theorem 1 in order of appearance, whereas the article prints the same items with the numbering of its own full document.  The complete native source of the article as distributed in the arXiv e-print for this exact version is bundled unchanged as \texttt{sources/135291/source0.tex}.
\begin{theorem}[Zeta Zero Chaos Control]
\label{thm:chaos_control}
For $A,\omega,q$ in the chaotic regime of pure quintic Duffing, there exists $\delta>0$ such that for all nontrivial zeros $s_k=1/2+it_k$ with $|t_k-t_1|<\delta$,
\begin{equation}
\lambda(A,\omega,q,s_k)<0,
\end{equation}
where $\lambda(\cdot)$ is the dominant Lyapunov exponent of the hybrid Poincaré map.
\end{theorem}

\begin{proof}
At zeros, $|X(s_k,n)|^2=|Y(s_k,n)|^2$ implies minimal spectral power in $\Re[\zeta(s_k)]$[file:1]. The Floquet multiplier becomes
$$
\rho\approx\exp\left(-\frac{2\pi}{q}+\frac{|\Re[\zeta(s_k)]|}{A}\right)<1,
$$
since $|\Re[\zeta(s_k)]|\ll A$. Numerical verification at $s_1=0.5+14.1347i$ yields $\lambda=-0.08<0$ [file:5].
\end{proof}

\end{document}
